\clearpage\begingroup\NoSourceHead\SourceLeaf{502b}
\renewcommand{\thefootnote}{\arabic{footnote})}\setcounter{footnote}{65}
\ArithmeticHead{45}{R}{WIRKUNGEN DER GAUSSSCHEN ARBEIT.}
\noindent\hspace*{\parindent}24. Am Schluss seiner Berechnung von $\delta U$ (Art. 26) macht \textsc{Gauss} die
Bemerkung, dass der Wert von $\delta U$ leichter durch geometrische Betrachtungen
hätte erhalten werden können. Eine solche mehr geometrische Lösung des
\textsc{Gauss}schen Variationsproblems hat J. \textsc{Bertrand} (1848) geliefert\footnote{\textit{Mémoire sur la théorie des phénomènes capillaires}, Journal de Mathématiques, Bd. 13 (1848),
S. 196—198; vgl. auch J. \textsc{Bertrand}, \textit{Sur les surfaces isothermes orthogonales}, ebenda Bd. 9 (1844), S. 119.}. Er denkt
sich den Übergang von der Fläche $U$ zur variierten Fläche $\overline U$ dadurch bewerkstelligt,
dass jeder Punkt von $U$ auf der durch ihn gehenden Normale
um ein unendlich kleines Stück $\delta n$ verschoben wird. Dadurch erhält man
die Fläche $\overline U$ bis auf eine unendlich schmale Zone zwischen der Gefässwand
und der von den Normalen entlang der Begrenzungskurve $P$ gebildeten geradlinigen
Fläche. Der Inhalt dieser Zone sei $\delta_2 U$, der übrige Teil von $\overline U$ sei
$\overline U_1$. Um nun zunächst den Zuwachs $\delta_1 U=\overline U_1-U$ zu berechnen, zerlegt
\textsc{Bertrand} die Fläche $U$ durch die beiden Scharen ihrer Krümmungslinien in
Elemente, und folgert dann aus den Eigenschaften der Krümmungslinien die
Proportion
\begin{equation*}\tag{63}
\frac{d\overline U}{dU}=\frac{(R+\delta n)(R'+\delta n)}{RR'},
\end{equation*}
woraus sich
\begin{equation*}\tag{64}
\delta_1 U=\int dU\mathpunct{.}\delta n\mathpunct{.}\left(\frac{1}{R}+\frac{1}{R'}\right)
\end{equation*}
ergibt. Zur Berechnung von $\delta_2 U$ wird die oben näher definierte Zone der
\clearpage\SourceLeaf{503}
\ArithmeticHead{46}{L}{BOLZA, GAUSS UND DIE VARIATIONSRECHNUNG.}
\noindent Fläche $\overline U$ durch Ebenen senkrecht zur Begrenzungskurve $P$ in Elemente zerlegt.
Eine einfache geometrische Infinitesimalbetrachtung liefert dann
\[
\delta_2 U=-\int dP\mathpunct{.}\delta n\mathpunct{.}\operatorname{cotg} i.
\]

Man erhält so zugleich eine geometrische Deutung für jedes der beiden
Integrale, aus denen sich der \textsc{Gauss}sche Ausdruck (45) für $\delta U$ zusammensetzt.
Durch dieselben Zerlegungen werden auch die Variationen der übrigen Glieder
von $\delta W$ und $\delta s$ berechnet.

Eine etwas andere Ableitung der Gleichung (64) gibt E. \textsc{Lamarle}\footnote{\textit{Exposé géométrique du calcul différentiel et intégral}, Mémoires couronnés et autres mémoires
publiés par l’Académie de Belgique, Bd. XV (1863), S. 572—573. \textsc{Lamarle} erwähnt die \textsc{Gauss}sche Formel (45) nicht.} (1863).
Auch er verschiebt jeden Flächenpunkt entlang der Normale, aber er zerlegt
das Flächenstück nicht durch die beiden Scharen seiner Krümmungslinien,
sondern durch irgend zwei zu einander orthogonale Kurvenscharen in Elemente
und berechnet die Variation des Flächeninhalts eines Elementes mit den
Seiten $d\sigma$, $d\sigma'$ aus der Formel
\[
\delta dU=d\sigma\mathpunct{.}\delta d\sigma'+d\sigma'\mathpunct{.}\delta d\sigma,
\]
indem er auf $\delta d\sigma$ und $\delta d\sigma'$ die allgemeine Formel für die Variation eines
Kurvenelementes bei normaler Verschiebung anwendet. — Nahe verwandt mit
diesen Betrachtungen von \textsc{Bertrand} und \textsc{Lamarle} ist die Ableitung der \textsc{Gauss}schen
Formel für $\delta U$, die H. \textsc{Weber}\footnote{\textsc{Weber}, loc. cit. 51), S. 69—70.} (1903) in den Anmerkungen zur deutschen
Ausgabe der \textsc{Gauss}schen Arbeit gegeben hat. Er weicht von \textsc{Bertrand}
hauptsächlich darin ab, dass er die Fläche $U$ nicht durch die beiden Scharen
ihrer Krümmungslinien oder sonst auf eine spezielle Art zerlegt, sondern die
Proportion (63) aus der Beziehung zwischen einem Flächenelement $dU$ und
dessen sphärischem Bild $d\omega$ ableitet: $dU=RR'd\omega$.

25. Am Schluss seiner Berechnung von $\delta U$ (Art. 26) hebt \textsc{Gauss} ferner
hervor, dass bei der ganzen Ableitung stillschweigend vorausgesetzt worden
war, dass auf der ganzen Fläche $U$ die Ordinate $z$ eine eindeutige Funktion
von $x,y$ ist, und dass $\zeta$ beständig positiv ist. Er zeigt, wie man sich durch
geeignete Zerlegung der Fläche von dieser beschränkenden Annahme befreien
\clearpage\SourceLeaf{504}
\ArithmeticHead{47}{R}{WIRKUNGEN DER GAUSSSCHEN ARBEIT.}
\noindent kann, und fügt hinzu, dass man denselben Grad der Allgemeinheit gleich von
vornherein hätte erreichen können, wenn man eine etwas andere Methode
angewandt hätte, womit er, wie schon H. \textsc{Weber} bemerkt, ohne Zweifel die
Behandlung des ganzen Problems in Parameterdarstellung gemeint hat. Für
diese Form der Darstellung hat H. \textsc{Weber} (1903) in den Anmerkungen zur
deutschen Ausgabe der \textsc{Gauss}schen Arbeit\footnote{\textsc{Weber}, loc. cit. 51), S. 63—69.} die Berechnung von $\delta U$ durchgeführt.
Aber schon vorher (1900) hatte \textsc{Kneser} in seinem \textit{Lehrbuch der
Variationsrechnung} (§§ 64, 65) die allgemeine Theorie der Extrema von Doppelintegralen
in Parameterdarstellung auf das \textsc{Gauss}sche Variationsproblem angewandt
und damit wohl die an Allgemeinheit, Eleganz und Strenge vollendetste
Lösung des Problems gegeben.

Bei dieser, wie es scheint, zuerst von \textsc{Weierstrass} in seinen Vorlesungen
angewandten Behandlungsweise werden die zulässigen Flächen in Parameterdarstellung
angenommen, wie dies schon \textsc{Poisson} getan hatte. Aber im Gegensatz
zu \textsc{Poisson}, dem die Parameterdarstellung nur ein vorübergehendes Hilfsmittel
gewesen war, dessen er sich so bald als möglich wieder entledigte, wird
hier an der Parameterdarstellung von Anfang bis zu Ende festgehalten. Dementsprechend
wird das zu untersuchende Doppelintegral von vornherein in
der Form
\[
J=\iint\Phi(x,y,z;\ x_u,y_u,z_u;\ x_v,y_v,z_v)\,du\,dv
\]
angenommen, wobei die Funktion $\Phi$ eine gewisse, zuerst von G. \textsc{Kobb}\footnote{\textit{Sur les maxima et les minima des intégrales doubles}, Acta Mathematica, Bd. 16 (1892), S. 68.} (1892)
aufgestellte Bedingung erfüllen muss, damit der Wert des Integrals nur von
der Fläche, nicht aber von der zufälligen Wahl des Parameters abhängt.
Variiert werden nur die Funktionen $x,y,z$, nicht aber die unabhängigen
Variabeln $u,v$.

Bei dem Integral $U$ ist dann $\Phi=\sqrt{EG-F^2}$ zu setzen, woraus sich ohne
Schwierigkeiten die \textsc{Gauss}sche Formel (45) für $\delta U$ ergibt. Dieselbe Formel
wendet \textsc{Kneser} dann mutatis mutandis auch zur Berechnung von $\delta T$ an, wie
dies auch schon \textsc{Pagani} getan hatte. Dabei ist nur zu beachten, dass hier
$\cos(4,5)=0$, und dass an Stelle der Richtung 7 jetzt die der Richtung 8
entgegengesetzte Richtung tritt. Die dreifachen Integrale $s$ und $\int z\,ds$ führt
\clearpage\SourceLeaf{505}
\ArithmeticHead{48}{L}{BOLZA, GAUSS UND DIE VARIATIONSRECHNUNG.}
\noindent \textsc{Kneser} mittels der \textsc{Gauss}schen Formel (42) auf Oberflächenintegrale zurück,
die über die gesamte Oberfläche der Flüssigkeit, also über $U$ und $T$ zu erstrecken
sind. So erhält er, wenn $d\sigma$ ein Element dieser Gesamtoberfläche
ist, für $s$ die im wesentlichen schon von \textsc{Gauss} gegebene Formel
\[
s=\tfrac13\int[x\cos(n,x)+y\cos(n,y)+z\cos(n,z)]\,d\sigma
\]
und weiter
\[
\int z\,ds=\int\frac{z^2}{2}\cos(n,z)\,d\sigma.
\]

Für die schliessliche Ableitung der Randbedingung bedient sich \textsc{Kneser}
eines der \textsc{Lagrange}schen Multiplikatorenmethode nachgebildeten Verfahrens,
das auch schon \textsc{Pagani} zu demselben Zweck angewandt hatte.
\par\clearpage\endgroup
